\section{A complete outgoing-angle capacity bound for physical classes}
\label{tc:section}

The cross-curve estimate of Section~\ref{sec:xd-cross-curve} can be
summed over all physical classes that meet one prescribed transversal.
Here we construct a genuine family of transversals which covers the
entire collision cylinder, including its coordinate seam, and perform
that sum.  Short classes and disconnected classes are included.  The
result is an unconditional bound for a geometric capacity moment and
for the unweighted sum over classes of a prescribed positive capacity.
The capacity is the angular size of the set of actual outgoing
directions present in a class; it is not a physical roof interval or
a replacement for the original same-roof comparison pool.

Throughout this section $R\in[R_-,R_+]=[9/20,47/100]$, the speed is
one, and $T_R$ is the original physical collision map.  Put
\[
 \kappa_R=R^{-1},\qquad a={100\over47},\qquad
 U={170\over9},\qquad \Gamma=\sqrt{1+(20/9)^2}.
\]
All curve lengths and tangent Jacobians are Euclidean arclength in
the physical coordinates $(s,\varphi)$, with $s=R\alpha$ and
$p=\sin\varphi$.  Fix once and for all the original $q=2$
homogeneity partition, its $k_0$, the curvature allowance $B_{\rm cv}$,
the short length $\ell_*$ and the constants $\theta_*<1,C_g$ of
Theorems~\ref{thm:uc-expansion} and~\ref{thm:uc-growth}.  Further
fixed shortening of $\ell_*$, if needed to apply their preliminary
short-graph arguments, is absorbed in that same curve-growth theorem.
None of these choices depends on the collision count, a record label,
a roof interval, or an outer parameter of the local-limit problem.

\subsection{The initially uncut affine family}

The actual outgoing direction is the circle-valued function
\begin{equation}\label{tc:direction}
 \pi_R(s,\varphi)=s/R+\varphi\pmod {2\pi}
       \quad\hbox{on }M_R=(\mathbb R/2\pi R\mathbb Z)
                                      \times(-\pi/2,\pi/2).
\end{equation}
In particular $\pi_R$ is globally defined at the collision-circle
seam.  For $\vartheta\in\mathbb R/2\pi\mathbb Z$, its fiber
$\Gamma_{R,\vartheta}$ is parametrized by
\begin{equation}\label{tc:fibers}
 -\pi/2<\varphi<\pi/2,\qquad
 s=R(\vartheta-\varphi)\pmod {2\pi R}.
\end{equation}
It is an injective physical stable-cone arc: in a lift its slope is
$d\varphi/ds=-\kappa_R$, its curvature is zero, and
\begin{equation}\label{tc:fiber-length}
 |\Gamma_{R,\vartheta}|=\pi\sqrt{1+R^2}
      \le S_{\rm aff}:=\pi\sqrt{1+R_+^2}.
\end{equation}
These are the collision states of parallel outgoing rays; no orbit
has been changed by introducing the foliation.  A fiber generally
crosses infinitely many current homogeneity strips.  We do not
partition it at those current strips before applying the next lemma.

Set
\begin{equation}\label{tc:affine-constants}
 S_{\rm im}:=\pi\sqrt{1+a^{-2}},\qquad
 N_{\rm im}:=1+\left\lceil {2S_{\rm im}\over\ell_*}\right\rceil,
 \qquad
 C_{\rm aff}:=\max\{1,N_{\rm im}\Gamma\theta_*C_g\}.
\end{equation}

\begin{lemma}[Growth from a nonhomogeneous affine stable source]
\label{tc:affine-growth}
Let $W$ be any open subarc of a fiber~\eqref{tc:fibers}, with
$|W|\le\ell_*$.  For every $N\ge1$, let
$\mathcal G_N^{\rm aff}(W)$ be obtained by $N$ actual backward
collisions, cutting at the actual regularity boundaries and at the
homogeneity boundaries at times $-1,\ldots,-N$, and making the fixed
short-length subdivisions after each collision.  Then
\begin{equation}\label{tc:affine-growth-bound}
 \sum_{V\in\mathcal G_N^{\rm aff}(W)}
                  \sup_V J_VT_R^N\le C_{\rm aff}.
\end{equation}
The pullbacks to $W$ of these open pieces are disjoint, and cover
arclength-almost all of $W$.  At every point of $W$ with an actual
regular backward record through time $-N$, the excluded cuts are
locally finite.  Such a point is either in one of the open pullbacks,
or is an endpoint of at least one of them.  In the latter case the
one-sided curve Jacobian has the value of the actual finite regular
derivative at that point.
\end{lemma}

\begin{proof}
Time reversal $I(s,\varphi)=(s,-\varphi)$ is an arclength isometry,
satisfies $IT_RI=T_R^{-1}$, and turns $W$ into a positive-cone
affine graph of slope $\kappa_R$ and curvature zero.  We first check
that the native one-step proof applies before this graph has been
made homogeneous.  The relative-angle and positive-foot estimates
in Section~\ref{sec:uc-complexity} require only the cone bounds and
the common short length.  They give at most $J_f$ actual first-hit
pieces, each meeting each target strip at most once.  The limiting
first-hit argument of Lemma~\ref{lem:uc-central}, including its
initial-grazing case, also uses no current-strip restriction.  There
is consequently at most one central target piece.  The exact
adapted expansion~\eqref{eq:uc-adapted} is valid for every positive
initial incidence, and its target-strip sum is therefore precisely
\eqref{eq:uc-onestep}.  Neither the initial incidence nor the current
strip index is divided out or summed in this argument.

The curvature formula~\eqref{eq:uc-second} and
Theorem~\ref{thm:uc-curvature} also use only the initial cone and
curvature bounds.  Every regular homogeneous first image hence
belongs to the inherited class with allowance $B_{\rm cv}$, even
though the source was not homogeneous.  Reversing the resulting
statement gives
\begin{equation}\label{tc:first-step}
 \sum_{V\in\mathcal H^-(W)}\sup_V J_VT_R
                         \le\Gamma\theta_*.
\end{equation}
Here $\mathcal H^-(W)$ consists of the maximal actual homogeneous
pieces of $T_R^{-1}W$, before their length subdivision.  The factor
$\Gamma$ is the one endpoint conversion from the adapted norm to
Euclidean arclength; it is not a factor at every subsequent collision.

On each of these graph pieces, $|d\varphi/ds|\ge a$ and
$\varphi$ is strictly monotone.  In a local lift,
\[
 d\ell=\sqrt{1+|d\varphi/ds|^{-2}}\,|d\varphi|
                    \le\sqrt{1+a^{-2}}\,|d\varphi|.
\]
Its total length is at most $S_{\rm im}$.  The graph cannot acquire
an extra turn by passing a coordinate seam: its lifted horizontal
length is at most $\pi/a<2\pi R_-$, while its collision angle stays
in $(-\pi/2,\pi/2)$.  A piece of length at most $S_{\rm im}$ can
be subdivided into at most $N_{\rm im}$ pieces of length at most
$\ell_*$, with the standard lower length $\ell_*/2$ when a long
piece is subdivided.  For example, a piece of length $L>\ell_*$
can be divided into $\lceil L/\ell_*\rceil$ equal parts.  Thus
the sum in~\eqref{tc:first-step}, after this fixed subdivision, is
at most $N_{\rm im}\Gamma\theta_*$.

Each new piece is now an actual homogeneous stable graph in
$I\mathcal W_R^u(B_{\rm cv})$.  The stable inverse form of
Theorem~\ref{thm:uc-growth} applies to its remaining $N-1$
backward collisions.  The chain rule, with the first-step supremum
paid only once per first piece, gives
\[
 \sum_{V\in\mathcal G_N^{\rm aff}(W)}\sup_V J_VT_R^N
 \le C_g\sum_{V_1\in\mathcal G_1^{\rm aff}(W)}
                                  \sup_{V_1}J_{V_1}T_R
 \le C_gN_{\rm im}\Gamma\theta_*.
\]
This includes $N=1$, since $C_g\ge1$.  At $N=0$ the analogous
sum is one.

For coverage, the physical first-hit cuts are finite on each short
graph, the target homogeneity cuts are countable, and the additional
length cuts are finite on each homogeneous image.  Finite iteration
of a countable collection remains countable.  Each regular branch
is injective, so the points removed on the original $W$ are countable.
This is the actual curve-null argument; it is not an inference from
the bulk measure of a singular set.

We need a little more than this null statement for endpoints.  Fix
a point $p\in W$ with a regular backward record of length $N$.
On a neighborhood of $p$ the finitely many actual iterates
$T_R^{-j}$ are smooth and have the same physical word.  Each of
their collision angles stays in a compact subinterval of
$(-\pi/2,\pi/2)$.  Only finitely many homogeneity levels can
therefore occur in a smaller neighborhood.  On the corresponding
stable image the derivative of its angle is nonzero, by cone
invariance.  A fixed such level cuts a one-sided regular graph at
most once.  The length subdivisions are locally finite there as
well: their interior cut points are separated by the fixed positive
image length $\ell_*/2$, and the finitely many image maps on this
neighborhood have finite derivatives.  Induction through the
finitely many times gives finitely many cut points near $p$.
Consequently, if $p$ was removed as a cut, a one-sided adjacent
open piece reaches it.  The original finite regular map and the
initial affine tangent are smooth at $p$, so their curve Jacobian
extends to that endpoint with the actual derivative value.  No
derivative at a true grazing state is used.
\end{proof}

\subsection{All actual words, strip faces, and unique lifts}

Let $D^-_{N,R}$ be the open set of collision states having an
actual regular backward record through $0,-1,\ldots,-N$.
All first collisions on these records are the physical first
collisions; an earlier unvisited tangency is part of the excluded
regularity boundary.  For $y\in D^-_{N,R}$, lift its current
scatterer to $O_0=0$.  Its past then determines the unique tuple
\begin{equation}\label{tc:labels}
 \bigl((O_0,O_{-1},\ldots,O_{-N});
                    h_{-1},\ldots,h_{-N}\bigr),
 \qquad T_R^{-j}y\in H_{h_{-j}}.
\end{equation}
The indices $h_{-j}$ use the fixed original disjoint assignment
of horizontal homogeneity boundaries.  In particular the terminal
past strip at $-N$ is included, and the current strip at $0$ is not.
Denote by $\mathfrak C_N(R)$ the nonempty sets of states with a
fixed tuple~\eqref{tc:labels}.  A class is the union of all its
physical components; it is not counted again for each disconnected
piece, section cut, or periodic lattice translate.  Boundary-only
faces allowed by the disjoint strip convention are also classes.

There are countably many such classes.  Indeed the uniform finite
flight bound gives a finite alphabet of physical center increments,
and a finite tuple of homogeneity indices is countable.  The
classes are disjoint and cover $D^-_{N,R}$ exactly.  Each is a
countable union of compact sets in the collision cylinder: a fixed
regular physical word has an open domain, the strip intervals with
their assigned endpoints are countable unions of closed sets, and
the finitely many relevant iterate maps are continuous on that
domain.  Intersect with a compact exhaustion of the domain to
obtain the assertion.  This observation also applies to classes
consisting only of artificial-boundary points.

Write $M_R'$ for the original set of points with regular two-sided
orbits, and let $j_N(x)=J^sT_R^N(x)$ be its true stable arclength
cocycle.  For $C\in\mathfrak C_N(R)$ define
\begin{equation}\label{tc:weight-capacity}
 d_N(C)=\sup_{y\in C\cap M_R'}j_N(T_R^{-N}y),\qquad
 E_R(C)=\pi_R(C),\qquad
 \operatorname{cap}_R(C)=|E_R(C)|.
\end{equation}
The supremum of an empty set is defined to be zero.  The measure
in the capacity is ordinary angle Lebesgue measure on
$\mathbb R/2\pi\mathbb Z$, of total mass $2\pi$.  Because the
projection is continuous and the classes are countable unions of
compact sets, $E_R(C)$ is Borel.  True stable invariance gives the
exact identity
\begin{equation}\label{tc:inverse-identity}
 j_N(T_R^{-N}y)
    ={1\over J_{E^s(y)}T_R^{-N}(y)}\qquad(y\in M_R').
\end{equation}
In particular the weights are finite, with the elementary bound
$d_N(C)\le\Gamma\Lambda_*^{-N}$.  There is no collision-flux,
section-conditioning, or radius normalization in this tangent
identity.

\begin{lemma}[One short test pays every class it meets]
\label{tc:one-test}
For an affine stable source $W$ as in Lemma~\ref{tc:affine-growth},
\begin{equation}\label{tc:one-test-bound}
 \sum_{\substack{C\in\mathfrak C_N(R)\\ C\cap W\ne\varnothing}}
                         d_N(C)\le3C_\times C_{\rm aff}
                         \qquad(N\ge1).
\end{equation}
This includes every disconnected class and every regular strip face
which meets $W$.
\end{lemma}

\begin{proof}
Fix a class with $d_N(C)>0$ which meets $W$, and choose a point
$p_C\in C\cap W$.  It need not have a regular future.  It does
have the required finite regular past.  For every $y\in C\cap M_R'$
the two backward words agree, and their angles lie in the closures
of the same strips at exactly $-1,\ldots,-N$.  The true stable
direction at $y$ and the affine tangent at $p_C$ are both in
$[-U,-a]$.  Corollary~\ref{cor:xd-backward-curves} and
\eqref{tc:inverse-identity} therefore give
\begin{equation}\label{tc:representative-weight}
 d_N(C)\le {C_\times\over J_WT_R^{-N}(p_C)}.
\end{equation}
Taking the class supremum has cost only this common constant.
No connecting path inside the class and no continuity of the true
stable field on a test curve have been assumed.

Let $\mathcal I_N(W)$ be the open pullbacks to $W$ of
$\mathcal G_N^{\rm aff}(W)$.  If $p_C$ is in an interval
$I\in\mathcal I_N(W)$, charge $C$ to that interval.  All interior
points of $I$ have the same actual word and the same strip indices,
so at most one distinct class can be charged to its interior.
If $p_C$ is a removed cut, the endpoint part of
Lemma~\ref{tc:affine-growth} allows it to be charged to an adjacent
interval.  An interval has at most two endpoints, and each endpoint
belongs to at most one class under the fixed disjoint convention.
It can therefore receive at most two additional distinct endpoint
classes.  This gives multiplicity at most three, regardless of the
number of components or intersections belonging to one class.
One representative is chosen per class, not per intersection.

If $V=T_R^{-N}I$ is its charged image piece, continuity at a
regular endpoint, when needed, implies
\[
 {1\over J_WT_R^{-N}(p_C)}\le\sup_V J_VT_R^N.
\]
Sum~\eqref{tc:representative-weight}, use multiplicity three,
and apply~\eqref{tc:affine-growth-bound}.  The argument applies
first to finite sets of classes and then to their increasing
countable union.  Classes of zero weight do not affect the sum.
\end{proof}

\subsection{Uniform capacity and the long-class consequence}

Put
\begin{equation}\label{tc:global-constants}
 N_{\rm aff}=2+\left\lceil {2S_{\rm aff}\over\ell_*}\right\rceil,
 \qquad K_{\rm tc}=3N_{\rm aff}C_\times C_{\rm aff},\qquad
 \beta_{\rm tc}={2a\over\sqrt{1+U^2}}>0.
\end{equation}

\begin{theorem}[All-class physical outgoing-angle capacity]
\label{tc:all-classes}
For every radius in the original family and every $N\ge1$,
\begin{align}
 \sum_{C\in\mathfrak C_N(R)}
       d_N(C)\boldsymbol 1_{E_R(C)}(\vartheta)
       &\le K_{\rm tc}
       &&\text{for every }\vartheta\in\mathbb R/2\pi\mathbb Z,
                                                    \label{tc:fiber-bound}\\
 \sum_{C\in\mathfrak C_N(R)}
       d_N(C)\operatorname{cap}_R(C)
       &\le2\pi K_{\rm tc}.                         \label{tc:capacity-bound}
\end{align}
Consequently, for every $\delta>0$,
\begin{equation}\label{tc:positive-capacity}
 \sum_{\substack{C\in\mathfrak C_N(R)\\
                 \operatorname{cap}_R(C)\ge\delta}}d_N(C)
                      \le{2\pi K_{\rm tc}\over\delta}.
\end{equation}
If $\mathfrak L_{N,R}(\delta)$ denotes the classes which contain
a connected physical unstable-cone arc of arclength at least
$\delta$, with slopes in $[\kappa_R,U]$, then
\begin{equation}\label{tc:long-class-bound}
 \sum_{C\in\mathfrak L_{N,R}(\delta)}d_N(C)
                      \le{2\pi K_{\rm tc}\over\beta_{\rm tc}\delta}.
\end{equation}
The constants are uniform before any supremum in radius or collision
count.  No connectedness, minimum area, or positive lower bound for
the capacity of an arbitrary class is required.
\end{theorem}

\begin{proof}
Parametrize $\Gamma_{R,\vartheta}$ by arclength on its open
interval of length $L\le S_{\rm aff}$.  For
$n=\lceil2L/\ell_*\rceil$, the open windows with centers at
$jL/n$, $0\le j\le n$, and half-width $L/n$, restricted to that
interval, cover it and each have length at most $\ell_*$.
There are at most $N_{\rm aff}$ windows.  They are the short
affine sources of Lemma~\ref{tc:one-test}.  They also cover a
point at the collision-circle seam, since the original fiber is
the smooth arc~\eqref{tc:fibers}, not two arcs stopped at that seam.

Every class with $\vartheta\in E_R(C)$ meets at least one window.
Sum~\eqref{tc:one-test-bound} over the windows; counting a class
more than once only increases this positive upper bound.  This
proves the pointwise assertion~\eqref{tc:fiber-bound}.  Apply
Tonelli's theorem to its countable nonnegative sum on the outgoing
angle circle to obtain~\eqref{tc:capacity-bound}.  Restricting to
classes of capacity at least $\delta$ proves
\eqref{tc:positive-capacity}.

For the last assertion orient a connected unstable-cone arc so
that $s$ and $\varphi$ increase in a lift, and write
$v=d\varphi/ds\in[\kappa_R,U]$.  Its outgoing-angle lift satisfies
\begin{equation}\label{tc:projection-derivative}
 {d\vartheta_{\rm lift}\over d\ell}
       ={v+\kappa_R\over\sqrt{1+v^2}}
                   \ge {2a\over\sqrt{1+U^2}}=\beta_{\rm tc}.
\end{equation}
Moreover,
\[
 0<d\vartheta_{\rm lift}
       =\left(1+{\kappa_R\over v}\right)d\varphi
       \le2\,d\varphi.
\]
Since the collision angle is in $(-\pi/2,\pi/2)$, the total
outgoing-angle variation is at most $2\pi$.  Thus passage to the
angle circle causes no loss of its one-dimensional image measure.
The projection of an arc of length at least $\delta$ has measure
at least $\beta_{\rm tc}\delta$ and is contained in $E_R(C)$.
Apply~\eqref{tc:positive-capacity} with this threshold.  For an
arc specified through a supremum of lengths, first take shorter
arcs and then pass monotonically to the stated length.  This proof
uses the actual monotone cone arc, and does not infer a long arc
from a large Euclidean diameter or a cell-shape assertion.
\end{proof}

\subsection{Exact pressure, optical, and observation scope}

Let $\mathcal P_N(R)$ denote the unchanged physical-word pressure
of Input~\ref{input:v112-geometric-pressure}.  Under $y=T_R^Nx$,
translation of the current lifted center to zero maps the forward
word $(0,O_1,\ldots,O_N)$ bijectively to the backward word
\[
 (0,O_{N-1}-O_N,\ldots,O_1-O_N,-O_N).
\]
This is reversal of the ordered record followed by one lattice
translation.  It preserves the actual collision states and tangent
norms.  Distinct periodic lifts of this same record do not make
new classes.  Since the homogeneity convention assigns every
finite regular state exactly once, including artificial strip faces,
\begin{equation}\label{tc:pressure-comparison}
 \mathcal P_N(R)
   =\sum_{w\in\mathcal W_N(R)}
        \sup_{\substack{C\in\mathfrak C_N(R)\\
                         C\text{ has reversed word }w}}d_N(C)
   \le\sum_{C\in\mathfrak C_N(R)}d_N(C).
\end{equation}
There is no deletion of boundary values in this comparison of
suprema.  All true grazing points remain outside $M_R'$, as in
the original definition of the stable derivative.  The finite
regular classes cover a full collision-measure set for each fixed
radius and finite time; a common null set for uncountably many
radii is neither used nor claimed.

Theorem~\ref{tc:all-classes} pays every term after multiplication
by its actual outgoing-angle capacity.  It also pays the entire
unweighted long-class sum~\eqref{tc:long-class-bound}, without a
planar-cell intersection or continuation assumption.  It does not
bound the unweighted sum on the right of
\eqref{tc:pressure-comparison} merely because small-capacity
classes are short.  More precisely, for any fixed $\delta>0$,
\begin{equation}\label{tc:remaining-pressure-interface}
 \mathcal P_N(R)\le {2\pi K_{\rm tc}\over\delta}
   +\sum_{w\in\mathcal W_N(R)}
       \sup_{\substack{C\in\mathfrak C_N(R)\\
                C\text{ has reversed word }w\\
                \operatorname{cap}_R(C)<\delta}}d_N(C).
\end{equation}
This states exactly the still-needed short-capacity payment along
this route; it is not an assertion that the last positive sum
vanishes or is uniformly bounded.  A complete physical renewal
or another genuine geometric payment must provide that additional
bound.  The geometric capacity in this section must in turn be
connected to the original conditioned physical-roof source before
it can supply a same-roof probability estimate.

The optical comparison already proved in module 313 remains
endpoint-exact:
\[
 {c(T_R^Nx)j_N(x)\over30000}
             \le\omega_N(x)\le30000c(T_R^Nx)j_N(x).
\]
For example the class weight
\[
 o_N(C):=\sup_{y\in C\cap M_R'}
                       {\omega_N(T_R^{-N}y)\over c(y)}
\]
satisfies $30000^{-1}d_N(C)\le o_N(C)\le30000d_N(C)$;
the corresponding capacity moment is at most
$30000\cdot2\pi K_{\rm tc}$.  This comparison retains the actual
terminal incidence, and introduces neither an incidence cutoff
nor a fictitious measure normalization.

For a fixed Borel observation event $A$ of the actual finite
record, one may replace $d_N(C)$ by the supremum over
$C\cap A\cap M_R'$ in every positive upper bound above.  Such
an event may specify the original section and return event, exact
displacement and collision labels, and any interval of the original
physical roof.  The upper bound is uniform in this restriction,
but gives no factor equal to the roof-interval width and does not
authorize summing separately repeated class suprema over arbitrary
additional labels.  In particular $1/(4\pi R)$ and $1/c_*$ remain
the inherited collision-flux and section-conditioning factors at
their original measure interfaces; they are absent from the tangent
and angle-capacity identities here.  Neither the original $m^2$
normalization nor the prescribed outer-$B$, collision-$\limsup$
order is changed.  The final raw pointwise local-limit and
four-coordinate inversion target consequently has precisely its
previous physical and probabilistic scope.

